\subsection*{An exact degree-six test for the genus-two route}

Let $L$ be the free Lie algebra over $\mathbf Q_\ell$ on
$x_1,y_1,x_2,y_2$, and write
\[
 z_1=[x_1,y_1],\quad z_2=[x_2,y_2],\quad
 c=[z_1,z_2],\quad a_1=[z_2,c],\quad a_2=-[z_1,c].
\]
The degree-six derivation
\[
 D(x_i)=[a_i,x_i],\qquad D(y_i)=[a_i,y_i]\quad(i=1,2)
\]
is nonzero and special. Indeed
\[
 D(z_1+z_2)=[a_1,z_1]+[a_2,z_2]
 =-[z_1,[z_2,c]]+[z_2,[z_1,c]]=-[c,c]=0.
\]
It is tangent to the conjugacy class of each of the four
generators by its defining formulas. The two $a_i$ have
different multidegrees and are nonzero. Their nonvanishing
can be seen in the free associative algebra: the commutators
$z_1,z_2$ freely generate a subalgebra because their least
words $x_1y_1,x_2y_2$ are two distinct equal-length code words,
so the least words of distinct products cannot cancel.

It is not an inner derivation. An equality $D=\operatorname{ad}w$
would make $w-a_1$ commute with both $x_1,y_1$, hence vanish
by their free-Lie centralizers; it would likewise give
$w=a_2$. This contradicts the different nonzero multidegrees.

For the diagonal CM torus with
\[
 x_i\mapsto\alpha_i x_i,\quad y_i\mapsto\beta_i y_i,
 \qquad \alpha_1\beta_1=\alpha_2\beta_2=t,
\]
both $z_i$ have character $t$, both $a_i$ have character
$t^3$, and conjugation on $D$ has character $t^3$.
For the $X_5$ automorphism, take
$(\alpha_1,\beta_1,\alpha_2,\beta_2)
=(\zeta_5,\zeta_5^{-1},\zeta_5^2,\zeta_5^{-2})$;
then $t=1$ and $D$ commutes with the order-five action.
Exchanging the two symplectic pairs exchanges $a_1,a_2$,
so this derivation also respects that symmetry.

This gives an exact nonzero vector of the relevant character
which is killed by each evaluation modulo $[x_i,L_6]$ and
$[y_i,L_6]$. Hence the rank-two determinant argument of
F.3.20 cannot be transferred by asserting that every
nonzero special genus-two derivation of character $\chi^3$
obstructs a generator's conjugacy class.

\paragraph{The exact map to the universal arithmetic quotient.}
Write $\omega=z_1+z_2$ and let $I=(\omega)$ be its Lie ideal.
The identity
\[
 c=[z_1,\omega]
\]
places $c,a_1,a_2$, and every $D(v)$ in $I$. Consequently the
induced derivation on the closed-surface Lie algebra $L/I$ is
zero. Filling the unique puncture quotients the actual
pro-$\ell$ group by the closed normal subgroup generated by
the boundary loop, whose initial term is $\omega$.
The natural graded map consequently factors through the
surjection $L/(\omega)\to\operatorname{gr}\pi_1(X_5)^{(\ell)}
\otimes\mathbf Q_\ell$. All values vanish through this
actual quotient map; injectivity of the graded surjection is
not needed.

For $r=0,1$, write $\Pi_r^{\mathrm{mod}}=\pi_1(M_{2,r})$,
$p_r:\Pi_r^{\mathrm{mod}}\to G_{\mathbf Q}$, and
$\Phi_r:\Pi_r^{\mathrm{mod}}\to\Gamma_{2,r}$ for the
universal outer monodromy representation. Put
\[
 F_r(m)=\Phi_r^{-1}(\Gamma_{2,r}(m)),\qquad
 \mathbf Q_{2,r}(m)=\overline{\mathbf Q}^{\,p_r(F_r(m))}.
\]
These are the exact field-tower definitions of Takao,
\emph{Braid monodromies on proper curves and pro-$\ell$
Galois representations}, RIMS~1709 (November 2010), pp.~4--5.
His Theorem~3.5, p.~21, proves
\[
 \mathbf Q_{2,0}(m)=\mathbf Q_{2,1}(m)\quad(m\ge1).
\]
It applies to genus~2 and includes filling the last puncture.
The older NTU~1994 Theorem~2.4, p.~115, assumes at least one
puncture and was not used to make that extra step.

Forgetting the marked point gives the actual morphism
$M_{2,1}\to M_{2,0}$; filling it in the universal curve gives
the compatible outer action on the quotient surface group.
The weight filtrations are compatible, so these maps induce
\[
\begin{array}{ccc}
 \operatorname{gr}_6\Gamma^{\mathrm{univ}}_{2,1}
 &\longrightarrow&
 \operatorname{gr}_6\Gamma^{\mathrm{univ}}_{2,0}\\
 q_{6,1}\downarrow&&\downarrow q_{6,0}\\
 \operatorname{Gal}(\mathbf Q_{2,1}(7)/\mathbf Q_{2,1}(6))
       \otimes\mathbf Q_\ell
 &\xrightarrow{\ =\ }&
 \operatorname{Gal}(\mathbf Q_{2,0}(7)/\mathbf Q_{2,0}(6))
       \otimes\mathbf Q_\ell.
\end{array}
\]
All graded spaces in this display are rational. The lower
identification is literally induced by the two equalities of
fields. Commutativity follows by applying $p_r$ to an element
of $F_1(6)$; the two arithmetic projections are the same.
To check descent to a graded monodromy image, an element with
zero image in degree~6 lies in $F_r(7)$ and therefore acts
trivially on $\mathbf Q_{2,r}(7)$. For rational classes one
first clears $\ell$-denominators; any torsion in an integral
graded quotient is then killed by a further power of $\ell$.
Thus the diagram uses actual group maps before tensoring.
The map $q_{6,1}$ is precisely the map in Nakamura~1996,
(4.7.1), p.~117.

Let $[D]$ denote the outer derivation class and put
\[
 J=\mathbf Q_\ell[D]\cap
     \operatorname{gr}_6\Gamma^{\mathrm{univ}}_{2,1}.
\]
The explicit closed-surface calculation gives zero on the
top horizontal image of $J$. The diagram therefore proves
the unconditional inclusion
\[
 \boxed{\ q_{6,1}(J)=0.\ }
\]
The displayed tangent vector cannot be a representative of
the nonzero arithmetic quotient used in F.3.20. Its membership
in the universal image itself has not been established; the
intersection statement does not assert that membership.
No procyclic arithmetic-normalizer subgroup is constructed.

\paragraph{Remaining calculation.}
The next task is to determine the image under $q_{6,1}$ of
the entire degree-six CM-weight tangent subspace that arises
from conjugacy equations, rather than this one explicit
line. The result above disposes of the line by a proved
forgetful morphism. It gives no assertion for the whole
subspace or for the split-prime genus-two normalizer problem.
